\documentclass[conference]{IEEEtran}
\IEEEoverridecommandlockouts
\usepackage{cite}
\usepackage{amsmath}
\usepackage{graphicx}
\usepackage{textcomp}
\usepackage{xcolor}
\usepackage{booktabs}
\usepackage{url}
\usepackage{array}
\def\BibTeX{{\rm B\kern-.05em{\sc i\kern-.025em b}\kern-.08em T\kern-.1667em\lower.7ex\hbox{E}\kern-.125emX}}

\begin{document}

\title{Meeting Real Time for On-Device Sign Language\\Translation on a 15\,W Jetson Orin Nano}

\author{\IEEEauthorblockN{Fnu Atisri}
\IEEEauthorblockA{\textit{The University of Texas at Dallas}\\
CS 6301: Low-Power Computer Vision}
\and
\IEEEauthorblockN{Tushar Goyal}
\IEEEauthorblockA{\textit{The University of Texas at Dallas}\\
CS 6301: Low-Power Computer Vision}
}

\maketitle

\section{Problem Statement}

The vision task is continuous sign language translation, which maps a sequence of
video frames of a person signing to an English sentence, so the output is a sentence
and not a sign label. A whole-body pose extractor turns each frame into 133 keypoints,
and a graph encoder followed by an mT5 text decoder turns that keypoint sequence into
English.

The deployment scenario is an assistive device placed on a counter or desk, which
translates American Sign Language into English text during a face-to-face exchange
between a Deaf signer and a hearing person who does not sign. The input is video of
the signer, cropped to the signer bounding box at native resolution and at the clip's
own capture rate. The device must run on the board rather than send frames to a
server. These are private medical and transactional exchanges, and the places they
happen, such as clinic desks and pharmacy counters, frequently have poor connectivity;
a translator that stops working without a network is not an assistive device.

The constraint that binds is \textbf{latency}. The pipeline must consume one second of
signed video in at most one second of wall clock at 15\,W, that is a real-time factor
of at most 1.00x. The reason is that the exchange is turn-taking. If the real-time
factor exceeds 1.00x then the backlog grows with every sentence, and by the third or
fourth turn the English text arrives after the conversation has moved on. Our baseline
measurement in Section~II gives \textbf{1.064x}, so the baseline fails this
constraint.

The accuracy floor is a \textbf{BLEU-4 loss of at most 1.00 point} against our
baseline of 23.16, and therefore a floor of 22.16. We set it at 1.00 because the same
checkpoint scores 22.67 in the paper, 22.53 replicating the authors' code on Colab in
bf16, and 23.16 on our own fp32 loop. That 0.6 spread is framework, precision and
frame subsampling, not run-to-run variance; a tighter floor would reject a stage over
a difference the evaluation path produces by itself.

\section{Baseline Model}

The baseline is the configuration we would deploy if this project did not exist. The
pose extractor is RTMW-l-m at 256$\times$192 producing 133 COCO-WholeBody keypoints,
exported to ONNX from the OpenMMLab deployment package and built as a TensorRT FP16
engine. The translator is the released Uni-Sign pose-only checkpoint~\cite{unisign},
which pairs a spatio-temporal graph encoder with mT5-base~\cite{mt5}, run in PyTorch
at FP32. The evaluation set is the OpenASL~\cite{openasl} test split, 976 clips, and
the board clip is a 255-frame, 8.51\,s sentence.

FP16 is the reference precision for the pose stage, as it is for most vision models on
this board. It is not the reference for the translator, because we measured mT5 in
FP16 returning token~0 at every step at a log probability of $-10.169$, exactly
$-\ln(26078)$ and therefore a uniform distribution from saturated logits. The baseline
mixes the two precisions deliberately.

Table~\ref{tab:baseline} reports our own measurements on course board
\texttt{jetson-lpcv-03} in the \texttt{lpcv} container at 15\,W mode~0, GPU clock
logged at 612\,MHz, occupancy confirmed idle before and after. The pose row is 897
timed inferences (3 repeats of 299 frames) at batch~1. The sentence row is 3 timed
sentences after a warm-up rather than the 1000 of the assignment protocol, because one
inference here is a whole 8.51\,s sentence and 1000 would hold a shared board for over
two hours; we report mean and standard deviation and carry this deviation forward.

\begin{table}[t]
\caption{Baseline measured by us on the course Jetson Orin Nano,
\texttt{lpcv} container, 15\,W mode 0, GPU 612\,MHz, batch 1.}
\label{tab:baseline}
\centering
\scriptsize
\setlength{\tabcolsep}{3pt}
\begin{tabular}{@{}p{0.40\columnwidth}p{0.52\columnwidth}@{}}
\toprule
\textbf{Property} & \textbf{Value} \\
\midrule
Models / weights & RTMW-l-m (OpenMMLab SDK ONNX); Uni-Sign pose-only + mT5-base \\
Input (px, fps) & 644$\times$720 signer crop, 29.97 \\
Precision, runtime & pose FP16 TensorRT 10.11; LM FP32 PyTorch 2.8 \\
Accuracy (BLEU-4 / ROUGE-L) & 23.16 / 43.17, 976 clips \\
Pose stage (ms/frame) & 19.7 \\
Sentence latency, P50 (ms) & 9052.8 $\pm$ 335.6 \\
\textbf{Real-time factor (x)} & \textbf{1.064} \\
Energy, total / dyn. (J/sent.) & 50.37 / 14.68 \\
Peak GPU memory (MiB) & 2577 \\
\bottomrule
\end{tabular}
\end{table}

Compression is needed because a smaller off-the-shelf model does not satisfy both
requirements. The pose stage is 6640.7 of the 9052.8\,ms, so a cheaper extractor is
the obvious alternative and we tested it. Substituting RTMPose-x, the extractor the
Uni-Sign authors name, costs \textbf{7.13 BLEU-4}, 95\% CI $[-10.94, -3.07]$, because
the translator was trained on RTMW-family keypoints and another extractor shifts its
input distribution. That is seven times the accuracy floor, and within the family
RTMW-l-m 256$\times$192 is already the smallest variant with a verified ONNX export.
The extractor is at its floor, so the remaining latency must come from running it
fewer times and from compressing the translator.

\section{Compression Plan}

We apply three stages, with a measurement between each. The order departs
deliberately from the Week~8 default of prune, distill, quantize. We prune first
because the vocabulary slice halves peak memory, and only once both models are
co-resident can any later stage be measured in the configuration we deploy. We
quantize second, and only the pose stage, because the translator's numerics forbid it.
We reduce the input rate last, because its accuracy cost must be charged against the
already-compressed model. Distillation is excluded; a fourth stage would weaken the
three.

\textbf{Stage 1: structured vocabulary pruning of the translator.} OpenASL uses 26{,}025
of the 250{,}112 mT5 tokens, so we slice those rows from the shared embedding and the
output projection, removing 344\,M of 587.7\,M parameters without retraining.
\emph{Re-measure:} BLEU-4 and ROUGE-L on 976 clips with a paired bootstrap, then board
latency, energy and peak memory. \emph{Failure mode A:} the keep set is selected using
the test split, leaking evaluation data into the architecture; we detect it by
re-running the census over train and dev only and re-scoring both keep sets.
\emph{Failure mode B:} a smaller softmax shifts beam scores even where the greedy
argmax is unchanged, so accuracy moves for a reason unrelated to capacity; we count
how many of the 976 predictions change.

\textbf{Stage 2: FP16 quantization of the pose extractor.} We build the extractor as an
FP16 TensorRT engine. \emph{Re-measure:} keypoint agreement against the FP32 engine,
as the fraction of confident keypoints within 5\,px, then end-to-end BLEU-4.
\emph{Failure mode A:} one activation overflows the FP16 range and silently destroys
the head. On RTMPose-x we measured the
ScaleNorm sum of squares at 142{,}293 against the FP16 maximum of 65{,}504. We detect
it before building, with a range scan that exposes every intermediate tensor and runs
the graph on ONNX Runtime CPU. \emph{Failure mode B:} TensorRT enables TF32 by
default, so an engine built as FP32 is not FP32 and the gate compares against the
wrong reference; clearing the TF32 builder flag moved our gate from 1.55e-3 to 7.2e-6
with all 2416 confident keypoints at 0.0\,px. \emph{Failure mode C:} layers the
builder silently keeps above the requested precision, the Week~7 case; we dump
per-layer precisions from the built engine rather than trust the build flag.

\textbf{Stage 3: input frame-rate reduction, 29.97 to 24\,fps.} The pose stage charges
per frame, so fewer frames is the largest single lever on the binding constraint.
\emph{Re-measure:} paired bootstrap on 976 clips at 30, 24, 16, 12 and 8\,fps, plus
board latency and energy over the rate by beam-width grid. \emph{Failure mode A:} the
accuracy cost is real but hidden by an interval too wide to exclude it. We made this error once, reading a 16\,fps cost of $-0.15$,
interval $[-3.20, +2.89]$ at $n{=}30$, as free, where $n{=}976$ later gave $-1.33$,
$[-2.00, -0.64]$. We fix $n{=}976$ in advance and refuse any claim whose interval
half-width exceeds the accuracy floor. \emph{Failure mode B:} OpenASL is not one frame rate
(76.5\% near 30, 22.2\% near 24), so a hardcoded source rate subsamples the wrong
frames; we derive each clip's rate from frames divided by duration.

Two techniques are excluded with a named reason. \emph{Weight-only INT8 on the
translator} is numerically sound at 22.79 BLEU-4 and a 4.05x smaller checkpoint, so it
does not inherit the FP16 failure. It is excluded because the decoder is bound by its
host-side step loop, so INT8 buys memory and not latency, and the int8 runtime is 10x
slower per token because 217 layers dequantize on every step: a number that improves
on paper with no change in measured latency. \emph{A TensorRT decoder engine set} is
token-identical at 19\,ms per token against 55 for PyTorch, but is not faster end to
end for the same reason and needs 2.07\,GB resident, which does not fit beside the
pose engine.

\textbf{Feasibility.} Fine-tuning is needed only if Stage~3 goes past 24\,fps. It runs
on Google Colab, to which we have access; one epoch over 20{,}000 clips with mT5
frozen trains 5.35\,M of 243.5\,M parameters in a few hours. We assume no desktop GPU
and train nothing on the board. We estimate 6 to 8 board-hours for the remaining
measurements, coordinated on the Teams channel; our 528-clip pose pass took 57.5
minutes. Beyond the \texttt{lpcv} container we need PyTorch~2.8,
\texttt{transformers} 4.44 and \texttt{onnxruntime}. All three stages should complete
by Milestone~2; the riskiest is Stage~3 at 16\,fps, where the cost exceeds the floor
and would need adaptation, and the fallback is 24\,fps, which needs none. Baseline and
Stage~1--2 board measurements are already taken.

\section{Metrics}

\label{m1}\textbf{Metrics.} The primary metric is the binding constraint of Section~I: end-to-end
sentence latency on the board at 15\,W, as P50 and P99 in ms and as a real-time factor
against the clip duration. Secondary metrics are energy per sentence in J, total and
dynamic, peak GPU memory in MiB, and per-stage energy integrated from each stage's own
power window rather than apportioned by latency share, since the pose stage draws
5.0\,W and the language model 6.2\,W. Accuracy is BLEU-4 and ROUGE-L on the 976-clip
test split. Parameter counts and checkpoint sizes are context, not results.

\label{m2}\textbf{Baseline.} Every comparison is against the Table~\ref{tab:baseline}
configuration, in the same runtime, on the same board and clip, and for accuracy on
the same 976 clips.

\label{m3}\textbf{Decision rule.} The project succeeds if the compressed pipeline reaches a
real-time factor of at most 1.00x at 15\,W \emph{and} the BLEU-4 loss against baseline
has a 95\% confidence interval whose lower bound is above $-1.00$ point. It fails
otherwise. The rule is stated on the interval and not on the point estimate, so a
result cannot pass on a loss that the data cannot actually bound.

\label{evidencestart}\textbf{Evidence.} Four ablations isolate each stage: vocabulary pruning on and off;
frame rate at 30, 24, 16, 12 and 8\,fps; beam width at 4, 2 and 1; and pose FP16
against FP32. We will plot accuracy against energy and against latency across the
stages. Every board row is 3 repeats with the DVFS state logged, and every accuracy
difference is a paired clip-level bootstrap on one shared set of draws.

\label{refsstart}
\begin{thebibliography}{00}
\bibitem{unisign} Z. Li, W. Zhou, W. Zhao, K. Wu, H. Hu, and H. Li, ``Uni-Sign: Toward
unified sign language understanding at scale,'' in \emph{Proc. ICLR}, 2025.
\bibitem{openasl} D. Shi, D. Brentari, G. Shakhnarovich, and K. Livescu, ``Open-domain
sign language translation learned from online video,'' in \emph{Proc. EMNLP}, 2022.
\bibitem{rtmw} T. Jiang, X. Xie, and Y. Li, ``RTMW: Real-time multi-person 2D and 3D
whole-body pose estimation,'' arXiv:2407.08634, 2024.
\bibitem{mt5} L. Xue, N. Constant, A. Roberts, M. Kale, R. Al-Rfou, A. Siddhant,
A. Barua, and C. Raffel, ``mT5: A massively multilingual pre-trained text-to-text
transformer,'' in \emph{Proc. NAACL}, 2021.
\bibitem{mmpose} MMPose Contributors, ``OpenMMLab pose estimation toolbox and
benchmark,'' \url{https://github.com/open-mmlab/mmpose}, 2020. Source of the RTMW-l-m
pretrained weights and ONNX export.
\bibitem{trt} NVIDIA, ``TensorRT 10.11 developer guide,'' 2025.
\end{thebibliography}

\section*{Contribution Statement}\label{contribstart}

The work is split by model stage, so that each partner owns a contract the other can
test against, and the interface between them is a fixed keypoint tensor format.
\textbf{Tushar Goyal} owns everything up to the keypoints: camera and frame handling,
the RTMW ONNX export and TensorRT engines, Stage~2 and its range scan, and all board
latency, power and energy measurement. \textbf{Fnu Atisri} owns everything after the
keypoints: the graph encoder and mT5 translator, Stage~1 vocabulary pruning, the
accuracy evaluation and the paired bootstrap, and the Stage~3 frame-rate study.
For this proposal, Atisri wrote Sections~I and~IV and the accuracy figures, Goyal
wrote Section~II's board measurements and Section~III's Stage~2, and both edited the
whole. The evaluation set and the measurement protocol are shared deliberately, so
that neither partner is the only person who measures the stage they built. Both
partners contributed equally.

\end{document}
